\section{Best response condition}

In the reduced-form environment, fix opponents' behavior $(\vminus,\aminus)$ and hence fix the
opponents' strength aggregator $\Sagg = \Sof{\vminus}{\aminus} > 0$ (Definition~\ref{def:Sof}).
Fix also $\delta>0$. Candidate $i$'s problem is then a binary choice between $a=0$ and $a=1$.

\subsection{Best response in difference form}

\begin{prop}[Reduced-form best response]\label{prop:br}
Fix $\Sagg >0$ and $\delta > 0$. A candidate with valence $v$ chooses $a = 1$ if and only if
\[
    \Deltainc{v}{\Sagg} \ge 0,
\]
where $\Deltainc{v}{\Sagg}$ is defined in Definition~\ref{def:Delta}.
\end{prop}

\begin{remark*}
This is a one-shot best-response condition: holding opponents fixed, the candidate takes the action
that maximizes the probability of ultimately winning office. Since the action set is $\{0,1\}$, it is
sufficient to compare expected utility under $a=1$ and $a=0$.
\end{remark*}

\begin{proof}
By Definition~\ref{def:Urf}, the reduced-form expected payoff from action $a$ is
$U_a(v;\Sagg,\delta) = \pnom{a}{v}{\Sagg}\Pge{v}{\delta}{a}$.
Thus $a=1$ is optimal if and only if
\[
    U_1(v;\Sagg,\delta) - U_0(v;\Sagg,\delta) \ge 0,
\]
which is exactly $\Deltainc{v}{\Sagg} \ge 0$ by Definition~\ref{def:Delta}.
\end{proof}

\subsection{Best response in ratio form (gain vs.\ cost)}

\begin{corollary}[Best-response as proportional gains and losses]\label{cor:BR}
Fix $\Sagg > 0$ and $\delta >0$. A candidate with valence $v$ chooses $a = 1$ if and only if
\[
    \Gfun{v}{\Sagg} \ge \Cfun{v}{\delta},
\]
where $\Gfun{v}{\Sagg}$ and $\Cfun{v}{\delta}$ are defined in
Definitions~\ref{def:Gbasic} and \ref{def:Cbasic}.
\end{corollary}

The ratio condition cleanly separates the forces in the candidate's decision:
$\Gfun{v}{\Sagg}$ is the proportional gain in nomination probability from taking $a=1$,
while $\Cfun{v}{\delta}$ is the proportional loss in general-election win probability.
The candidate takes $a=1$ exactly when the primary-stage gain outweighs the general-election cost.


\begin{proof}
Start from $\Deltainc{v}{\Sagg} \ge 0$, i.e.
\[
    \pOne{v}{\Sagg}\PhiStd(v-\delta) \ge \pZero{v}{\Sagg}\PhiStd(v).
\]
Since $\pZero{v}{\Sagg}>0$ and $\PhiStd(v-\delta)>0$, divide both sides by
$\pZero{v}{\Sagg}\PhiStd(v-\delta)$ to obtain
\[
    \frac{\pOne{v}{\Sagg}}{\pZero{v}{\Sagg}} \ge \frac{\PhiStd(v)}{\PhiStd(v-\delta)},
\]
which is $\Gfun{v}{\Sagg} \ge \Cfun{v}{\delta}$.
\end{proof}

\subsection{Best response in log form}

\begin{corollary}[Best-response in log form]\label{cor:BRlog}
Fix $\Sagg > 0$ and $\delta >0$. A candidate with valence $v$ chooses $a = 1$ if and only if
\[
    \logg{\Gfun{v}{\Sagg}} \ge \logg{\Cfun{v}{\delta}},
\]
i.e.
\[
    \logg{\pOne{v}{\Sagg}} - \logg{\pZero{v}{\Sagg}}
    \;\ge\;
    \logg{\PhiStd(v)} - \logg{\PhiStd(v-\delta)}.
\]
\end{corollary}


The log form is numerically convenient: it converts the multiplicative gain--cost comparison
into an additive inequality. This is more stable in simulations when probabilities are close
to $0$ or $1$.


\begin{proof}
By Corollary~\ref{cor:BR}, $a=1$ is optimal if and only if $\Gfun{v}{\Sagg} \ge \Cfun{v}{\delta}$.
Since $\Gfun{v}{\Sagg}>0$ and $\Cfun{v}{\delta}>0$, taking $\log$ of both sides yields the stated
inequality. Expanding $\Gfun{v}{\Sagg}$ and $\Cfun{v}{\delta}$ gives the final expression.
\end{proof}